\chapter{What the instruments could give}\label{ch:reach}

\section*{The place}
Part~\ref{part:3} built two device gauges and Part~\ref{part:4} built the cell instrument; Chapter~\ref{ch:onegauge} set the three on one gauge. They read
one quantity in three places: a cell's copy error against the healthy cell of its kind and the thermal floor, a qubit's gate error against the
thermal floor at the temperature its record sees, and a chip's switching energy against the Landauer floor. This chapter says what each could give if it is carried through, what
each has measured so far, and what has to happen before any of it is relied on. Every claim carries its status. Nothing here is commissioned
except where stated.

\section{What the three have in common, and what is new}

Each instrument reads \emph{one object} against two things that belong to that object: its own healthy or as-built state, and the floor
physics sets for it. No group of other objects enters the number. That is the difference from the methods each field uses now.
\begin{center}\small
\begin{tabular}{@{}p{0.2\textwidth}p{0.36\textwidth}p{0.36\textwidth}@{}}\toprule
 & common practice & these instruments \\\midrule
reference & a group: cases against controls, a training set, a curve built from many people & the same cell type healthy; the device as built; the physics floor \\
result & a probability, a class label, or a position in a distribution & a reading with its interval, on one gauge, 1 = healthy \\
outside the group it was built on & weakens; must be re-trained or re-calibrated & the reference does not move with the group \\
one subject over time & compared with a group at each time & compared with itself; the difference is the measurement \\
what it is not able to read & seldom stated & stated and printed; a reading it cannot make is refused \\\bottomrule
\end{tabular}
\end{center}
A group-based method answers \emph{how does this person compare with people like them}. These instruments answer \emph{how far has this
object moved from its own healthy state, and how far is it from the floor}. The two questions are complementary. The second has not had an
instrument in biology, because it needs a physical floor and a reference that is measured cell type by cell type rather than averaged over
people; IAM's Law supplies the floor, and the atlas supplies the reference.

\section{The cell}
The cell instrument is in development: one cell type (neutrophils) on one platform (EPIC v1 arrays), with IAM-A on single-molecule
sequencing beside it. What follows is what it could give, with what exists for each.

\paragraph{Cancer: reading the patient's own tissue against itself.} A tumour is a group of cells whose pattern has blurred. The potential is a reading for one patient that needs no reference group of
other patients: a biopsy read against the same patient's normal tissue, and later a blood draw read against the same person's earlier draw.
The blood reading, the size of a change it can see, and where breach and the cancer region lie on the gauge are to be measured. \openprob

\paragraph{One person over time.} A person can be measured again. Two draws months apart give a difference in which the person's baseline and
the laboratory's character cancel, which is the strongest form of the measurement and the one cosmology can never have with its single sky
(Chapter~\ref{ch:skytools}). Serial readings would turn the instrument from a snapshot into a trajectory: watching a high-risk person, a
treated patient after therapy, or a pre-cancerous condition, each against their own past. \prediction{} No serial specimens have yet been read
on the current chain.

\paragraph{Medicines and trials.} A drug that acts on DNA methylation can be read directly as a change in copy error. The potential is an endpoint that is per patient and mechanistic: whether a drug changed
the fidelity of the cells it targets, in each person, read against that person's own pre-treatment specimen, and on which channel (marks lost
or marks gained). No patient specimens have been read this way. \openprob

\paragraph{Ageing.} Ageing clocks~\cite{Horvath2013} place a person against an age curve built from many people. This instrument would instead read the
rate at which one person's cells lose their pattern, against their own earlier reading. Cultured lung fibroblasts~\cite{Cruickshanks2013} show what a channel-level
reading sees: senescent cells hold the methylated channel at the healthy rate and clean the unmethylated channel (0.69), while immortalised
cells carry more error on the methylated channel and less on the unmethylated one, and the two cancel in a combined number
(Chapter~\ref{ch:onegauge}). \measured{} Whether an intervention slows the loss in one person, read serially, is a test this design can make.
\prediction

\paragraph{Animals and temperature.} IAM's Law prices a bit at the local temperature. At a fixed holding energy the copy-error floor of a cell
at 10\,\textdegree C is 0.78 times the human floor $H(\varepsilon_0)$, and at a dog's 38.5\,\textdegree C it is 1.012 times. \calc{} That makes the
instrument usable on any vertebrate whose cell types can be referenced, and makes temperature itself a variable: farmed and wild fish at
water temperature, ectotherms under warming, veterinary ageing and disease in companion animals. \prediction

\paragraph{The laboratory itself.} Before any biology, the chain checks the specimen and the array: custody hashes, detection, call rate, sex,
conversion. A reading that states its instrument, its reference and its failures is
something a second laboratory can repeat.

\paragraph{A research tool for genetics.} The residual sky, the genome-distance spectrum and the C-score give a geneticist the scale at
which a departure lives, not only which sites moved (Chapter~\ref{ch:skytools}). \measured{} for the healthy case; untested on disease.

\section{The qubit}

The qubit gauge (Chapter~\ref{ch:ascoreqc}) reads a quantum processor's two-qubit gate error against the thermal floor of its record, at the
temperature the record actually sees (Chapter~\ref{ch:qplatforms}). What it could give engineers:
\begin{itemize}
\item \textbf{Where effort pays.} A gate error split into its parts, each set against the floor of the bath that drives it, says whether the
next gain comes from materials, from control, or from temperature, and how much room each channel has before its floor \interp.
\item \textbf{A floor from numbers the laboratory already has.} The gap, the qubit's own effective temperature, the gate time and $T_1$ give the
per-gate floor $p_{\rm eq}t_g/T_1$ with nothing assigned: $6.2\times10^{-7}$ for a 5~GHz transmon at 35~mK with $T_1=68~\mu$s and a 40~ns gate,
three orders of magnitude below present two-qubit errors \calc. Halving the qubit's own temperature lowers its held-record floor from
$1.05\times10^{-3}$ to $1.1\times10^{-6}$ \calc; whether the gate error follows is a measurement any refrigerator sweep can make on one device.
\openprob
\item \textbf{The quasiparticle limit.} An excess quasiparticle fraction of $10^{-7}$ alone limits an aluminium transmon at 5~GHz to
$T_1\approx0.24$~ms (Chapter~\ref{ch:xqp}). \calc{} The proposed internal source makes a test that existing devices can run: the background should
scale with the active-fluctuator count over island volume, $N/(\tau_{\rm TLS}V)$. \prediction
\end{itemize}

\section{The chip}
Figure~\ref{fig:cg_floor_Tj} (Chapter~\ref{ch:chipgen}) shows where the Landauer floor sits on one chip's gauge and how it moves with the chip's own
junction temperature.

The chip gauge (Chapter~\ref{ch:cmos}) reads a processor's energy per transistor switch against the Landauer floor at its junction temperature. For
the AMD Ryzen 9 9950X, $E=\mathrm{TDP}/(Nf)=170\,\mathrm W/((20.0$--$20.6)\times10^9\times4.3\,\mathrm{GHz})=(1.92$--$1.98)\times10^{-18}$~J~\cite{AMD9950X}, 576--593 times
$\kB T\ln2$ at 75\,\textdegree C; the transistor count, from die-level reports, is the least certain input. \calc{} What it could give designers is one number that tracks how close a design is to the physical floor, across
the generations of a design line on one definition of each input, and a statement of how much room is left: a factor of 576--593, close to three orders of magnitude, for the 9950X. The figure is an average over the whole chip and its idle transistors;
a per-block reading needs the design's own power and activity data. \openprob

\section{What none of this does}

None of these instruments diagnoses, prescribes or replaces the person who decides. The cell instrument reports a state and refers; it names
no condition, and no claim that it can or cannot detect something is made until the chain has been run on that question under seal. None of the
cell applications above is available for use until the instrument is commissioned for that cell type, platform and specimen, one at a time.

\section{What it would take}
\begin{enumerate}
\item Commission neutrophils on EPIC arrays end to end; then add one cell type at a time, each with its own healthy reference.
\item Read serial specimens of the same people and measure the smallest change a difference can see.
\item Place breach and the cancer region on the gauge from per-person data: tumour/normal pairs, the drug series, blood cancers.
\item Derive the Met-A floor, and read the single-molecule drug response on patient specimens.
\item For fish, remove library duplication and conversion failure, then repeat the temperature test.
\item For qubits, publish the gap, the qubit's own temperature, the gate time, $T_1$ and the isolated two-qubit error of one device together, and test the
quasiparticle prediction; for chips, read per-block energy where makers publish it.
\end{enumerate}

Each item on this list is a defined project with a clear end point. A geneticist, a device physicist or a chip engineer can take one up and carry it to its result.
